\xhdr{KG-based RAG}
Knowledge graphs (KGs) are widely used as knowledge sources in retrieval-augmented generation (RAG) \cite{han2024retrieval,guo2025empowering，yuthought} systems to enhance large language models (LLMs) with both relational and textual information for answering complex queries \cite{edge2024local,dong2024advanced,cui2024prompt,xia2024knowledge}. 
A core challenge in KG-based RAG lies in retrieving relevant knowledge from KGs in response to a given query. 
Recent methods addressing this problem can be broadly categorized into two approaches. 
Retrieval-based methods \cite{yasunaga2021qa,mavromatis2024gnn,he2024g,qiu2025graph,zou2025weak} leverage pretrained language models to encode the textual content in KGs into embeddings and use small models, such as MLP and GNN, to retrieve relevant information. 
In contrast, agent-based methods \cite{sunthink,zhang2024knowgpt,ma2024think,luoreasoning,liu2025hoprag} employ LLMs as agents that iteratively traverse the KG to locate supporting evidence. 
While both paradigms have shown promise in knowledge graph question answering (KGQA) \cite{liretrieval}, 
their effectiveness in retrieving diverse and high-quality candidates for complex queries remains limited.
Furthermore, complex queries often require retrieval from text-rich KGs, necessitating the joint consideration of both relational structure and text content. 
Although recent studies \cite{ma2024think,lei2025mixture} have begun to explore this setting and tackle complex queries, enhancing the diversity and accuracy of KG-based RAG is still underexplored.

\xhdr{Process Reward Models}
Process Reward Models (PRMs) \cite{lightman2023let,lambert2024rewardbench} have shown great promise in guiding LLMs with process supervision and have been adopted in many domains such as complex reasoning \cite{cui2025process}, alignment \cite{ouyang2022training}, and planning \cite{choudhury2025process}.
The key to PRMs is to construct a preference dataset with process supervision \cite{peng2025rewarding}.
Previous works obtain the process supervision from human feedback and LLM evaluation.
However, fine-grained process-level supervision is expensive for KG-based RAG due to the potentially vast search space of KGs and the difficulty of accessing the intermediate state during retrieval.
Although early explorations are made to use PRM to guide the retrieval process of RAG on unstructured knowledge bases \cite{li2025search,shi2025searchrag,huang2025rag,xiong2025rag}, they still need a preference dataset with process-level supervision.
How to guide the retrieval process of KG-based RAG on structured KGs without process supervision data is still challenging.

\xhdr{GFlowNet}
GFlowNet \cite{bengio2023gflownet} aims to sample diverse and high-quality candidates from an unnormalized density and has received increasing attention in sampling from discrete and vast spaces \cite{ghari2023generative,lu2024cell,deleu2023joint,deleu2022bayesian,m2024gflownet}.
The goal of GFlowNet is to learn a policy that can lead to the terminal states with the likelihood in proportion to their rewards \cite{zhang2022generative}.
Some objects are proposed to optimize GFlowNet by regularizing the state flows and their transitions \cite{tongimproving,malkin2022trajectory,madan2023learning}.
Recently, GFlowNet has also been introduced to improve the generative performance of LLMs and diffusion models by promoting diversity in the decoding process \cite{venkatramanamortizing,huamortizing,yu2024flow,kang2025gflowvlm,lingamenhancing} .
Differently, our work focuses on aligning the retrieval results of KG-based RAG with the knowledge required for real-world queries by estimating the state flow in multi-step retrieval.
Moreover, we introduce a local exploration strategy to avoid visiting less-valued states, thus efficiently optimizing the detailed balance.